\subsection{The completed local ring of the rank-five circuit}
\label{nl:section}

The six normalization points at the rank-five circuit admit an exact
description of their gluing. We retain the matrices
$M,Q_1,\ldots,Q_6$ of~\eqref{eq:M} and~\eqref{eq:Q}, their span
$L_0\subset A_0\otimes B_0$ with $\dim A_0=\dim B_0=3$, and the tensor
$T_5=\sum_{i=1}^5Q_i\otimes f_i$ of~\eqref{eq:T-first}.
The subspaces and tensor are embedded by coordinate inclusions whenever
the ambient factor dimensions are larger.

\begin{theorem}\label{nl:theorem}
Let $a,b\geq3$, $a+b\geq7$ and $c\geq5$, and put
$X=X_5(a,b,c)$. At the specified padded tensor $T=T_5$, set
\begin{equation}\label{nl:parameters}
 q=2(a+b-6),\qquad h=3a+3b+5c+2.
\end{equation}
There is an isomorphism of complete local $\C$-algebras
\begin{equation}\label{nl:completed-ring}
 \widehat{\OO}_{X,T}\simeq
 \frac{\C[[z_1,\ldots,z_h,\,
              x_{i,j}:1\leq i\leq6,\ 1\leq j\leq q]]}
      {(x_{i,j}x_{k,\ell}:i\neq k)}.
\end{equation}
In particular,
\begin{align}
 \dim\OO_{X,T}&=h+q=5(a+b+c-2),\label{nl:dimension}\\
 \depth\OO_{X,T}&=h+1=3a+3b+5c+3,\label{nl:depth}\\
 \operatorname{edim}\OO_{X,T}&=h+6q,\qquad
 e(\OO_{X,T})=6.\label{nl:embedding-multiplicity}
\end{align}
Here $e$ denotes Hilbert--Samuel multiplicity. The normalization of the
completed ring is
\begin{equation}\label{nl:normalization}
 \prod_{i=1}^6\C[[z_1,\ldots,z_h,x_{i,1},\ldots,x_{i,q}]].
\end{equation}
Every two branches of the completed local ring meet the same reduced smooth
$h$-dimensional formal germ. The conductor into~\eqref{nl:normalization}
is generated by all the $x_{i,j}$, and its quotient is
$\C[[z_1,\ldots,z_h]]$.

The assertions are preserved by simultaneously permuting the tensor
factors and transporting this specified witness.
\end{theorem}

\begin{proof}
\emph{The core chart.}
By Lemma~\ref{norm:core}, the scheme
\[
 \PP L_0\cap\Seg(\PP A_0\times\PP B_0)
\]
consists of exactly the six reduced points $[Q_i]$. Every five of the
$Q_i$ are independent, and their only relation, up to scalar, is
$\sum_iQ_i=0$.
Write $W_0=A_0\otimes B_0$ and consider
\[
 G_0=\Gr(5,W_0),\qquad
 Z_0=\{(L,p)\in G_0\times\Seg(\PP A_0\times\PP B_0):
                         p\in\PP L\}.
\]
The Grassmannian $G_0$ is smooth of dimension $20$.
Projection of $Z_0$ to the Segre is a Grassmann bundle: over a line
$p\subset W_0$ its fiber is $\Gr(4,W_0/p)$, of dimension $16$.
Thus $Z_0$ is also smooth of dimension $20$.

The projection $Z_0\to G_0$ is proper, and its fiber at $L_0$ is the
six-point reduced scheme just described. Upper semicontinuity of fiber
dimension, followed by proper quasi-finiteness, makes the projection
finite after shrinking around $L_0$. At each point of that fiber the
relative cotangent space is zero, so the morphism is unramified there.
The source and target are smooth of the same dimension; its differential
is therefore an isomorphism, and the morphism is \'etale at those
points. Remove the proper image of the non-\'etale locus. After a further
shrinking it is finite \'etale of degree six.

After an \'etale base change we may label the six sections
$p_1,\ldots,p_6$. Work on a sufficiently small smooth \'etale
neighborhood $U_0$ of $L_0$, with a chosen complex point over $L_0$.
Trivialize the line bundles of the six sections and choose regular
nonzero representatives
\begin{equation}\label{nl:regular-representatives}
 q_i=a_i\otimes b_i\in L,\qquad \sum_{i=1}^6q_i=0.
\end{equation}
To obtain the last equality, start with arbitrary representatives.
Every five remain a basis near $L_0$, so the unique relation among the
six has all coefficients invertible. Absorb those coefficients into the
representatives and into one of their factors. Choosing a nonzero
matrix entry of each rank-one representative makes its factorization
into $a_i,b_i$ regular. After shrinking again, the $a_i$ span $A_0$
and the $b_i$ span $B_0$. The completed local ring of $U_0$ at the
chosen complex point is the same as that of $G_0$ at $L_0$.
Only this core incidence has been made \'etale; no flatness of the
larger-format incidence is being assumed.

\medskip
\emph{Graph coordinates for nearby five-planes.}
Choose complements
\[
 A=A_0\oplus U,\qquad B=B_0\oplus V,\qquad
 \alpha=\dim U=a-3,\quad \beta=\dim V=b-3.
\]
Near $L_0$, projection to $W_0$ identifies a five-plane with the graph
over its projected $L\in G_0$ of three maps
\begin{equation}\label{nl:graph-maps}
 F:L\longrightarrow U\otimes B_0,\qquad
 G:L\longrightarrow A_0\otimes V,\qquad
 K:L\longrightarrow U\otimes V.
\end{equation}
Pull back this open Grassmannian chart to $U_0$.
Projection is injective on each such graph, so a nonzero rank-one
point of the graph projects to a nonzero rank-one point of $L$.
Its projective image is one of the six $p_i$, and there is at most one
lifted line over each $p_i$.

The lift whose core component is $q_i=a_i\otimes b_i$ is rank one
exactly when there are $u_i\in U$, $v_i\in V$ with
\begin{equation}\label{nl:lift-equations}
 F(q_i)=u_i\otimes b_i,\qquad
 G(q_i)=a_i\otimes v_i,\qquad
 K(q_i)=u_i\otimes v_i.
\end{equation}
Indeed, expanding $(a_i+u_i)\otimes(b_i+v_i)$ gives these three
components. Conversely, a rank-one tensor with core component exactly
$a_i\otimes b_i$ has this factorization after reciprocal rescaling of
its two factors. This rescaling does not change the equations.

\medskip
\emph{Residues and their common zero locus.}
Consider the vector-bundle map
\begin{equation}\label{nl:residue-map}
 \begin{split}
 \Phi_A:\Hom(L,U\otimes B_0)&\longrightarrow
           \bigoplus_{i=1}^6 U\otimes(B_0/\langle b_i\rangle),\\
 F&\longmapsto
       \bigl(F(q_i)\bmod U\otimes\langle b_i\rangle\bigr)_{i=1}^6.
 \end{split}
\end{equation}
This map is surjective. Given arbitrary quotient classes, choose
representatives $w_i\in U\otimes B_0$. Since the $b_i$ span $B_0$,
there exist $u_i\in U$ such that
\[
 \sum_i u_i\otimes b_i=-\sum_iw_i.
\]
The six vectors $w_i+u_i\otimes b_i$ then sum to zero. By
\eqref{nl:regular-representatives} and the independence of every five
$q_i$, they are the values of a unique linear map on $L$. This also
proves surjectivity when $\alpha=0$.

The domain of $\Phi_A$ has rank $15\alpha$ and its target has rank
$12\alpha$, so its kernel has rank $3\alpha$. It consists precisely of
\begin{equation}\label{nl:kernel-A}
 F=(P\otimes\operatorname{id}_{B_0})|_L,
             \qquad P\in\Hom(A_0,U).
\end{equation}
These maps lie in the kernel. The map $P\mapsto F$ is injective because
the $a_i$ span $A_0$; equality with the kernel follows from the ranks.
The same construction for $G$ gives a surjection
\[
 \Phi_B:\Hom(L,A_0\otimes V)\longrightarrow
        \bigoplus_{i=1}^6(A_0/\langle a_i\rangle)\otimes V
\]
whose kernel is exactly
\begin{equation}\label{nl:kernel-B}
 G=(\operatorname{id}_{A_0}\otimes Q)|_L,
             \qquad Q\in\Hom(B_0,V).
\end{equation}

After shrinking $U_0$, split these surjections of vector bundles and
choose frames for their rank-two quotient factors. This gives coordinates
$P,Q$ and six groups
\[
 x_i=(x_{i,1},\ldots,x_{i,q}),\qquad
 q=2\alpha+2\beta.
\]
The first two equations of~\eqref{nl:lift-equations} hold for every
index other than $i$ if and only if $x_k=0$ for every $k\neq i$.
If every residue vanishes, they hold for all six indices with
$u_i=P(a_i)$ and $v_i=Q(b_i)$. Applying $P\otimes Q$ to
\eqref{nl:regular-representatives} gives
$\sum_i u_i\otimes v_i=0$, so the six last equations are consistent
and determine the common value
\begin{equation}\label{nl:common-K}
 H_0=(P\otimes Q)|_L.
\end{equation}
This common locus is smooth of dimension $20+3\alpha+3\beta$.
The assertion follows from these exact family coordinates, rather than
only from tangent dimensions.

\medskip
\emph{Simultaneous straightening of the six configuration branches.}
Fix $i$. On the residue axis $x_k=0$ for $k\neq i$, require the five
points with indices $k\neq i$ to lift. Their $q_k$ form a basis of $L$.
The first two equations of~\eqref{nl:lift-equations} determine $u_k,v_k$
uniquely, and the last equations determine $K$ on that basis. Denote
this map by
\[
 K=H_i(L,P,Q,x_i).
\]
It is regular: the basis determinant and the selected nonzero factor
entries that must be inverted are units on the chosen neighborhood.
It is bilinear in the $F,G$ variables over the core chart. At $x_i=0$
it agrees with $H_0$. Thus the locus $B_i$ containing these five
lifted lines is an embedded smooth closed graph with coordinates
$L,P,Q,x_i$.

Extend each $H_i$ to all residue coordinates by ignoring $x_k$ for
$k\neq i$, and make the single change of coordinates
\begin{equation}\label{nl:simultaneous-translation}
 H=H_0+\sum_{i=1}^6(H_i-H_0),\qquad K'=K-H.
\end{equation}
This is invertible because $H$ depends only on the other graph
coordinates. On the $i$th residue axis every $H_k$ with $k\neq i$
equals $H_0$, so $H=H_i$. Consequently all six graphs become
\begin{equation}\label{nl:straight-branches}
 B_i=\{K'=0,\ x_k=0\text{ for every }k\neq i\}.
\end{equation}
This simultaneous description proves independence of all the normal
groups and their scheme-theoretic pair intersections. Pairwise
transversality alone is not needed to infer either conclusion.

\medskip
\emph{Exhaustion of limiting branches.}
Let $S=\Seg(\PP A\times\PP B)$, and let $\mathcal H$ be the reduced
closure in $\Gr(5,A\otimes B)$ of the spans of five independent points
of $S$. We next show that its reduced local germ is exactly
$\bigcup_{i=1}^6B_i$.

The point incidence
\[
 Z=\{(L,p):p\in\PP L\cap S\}
       \longrightarrow\Gr(5,A\otimes B)
\]
is proper. Its fiber at the padded $L_0$ is still exactly the six
reduced points, by Lemma~\ref{norm:core}: padding adds neither rank-one
directions nor tangent directions in that fiber. As for the core
incidence, proper quasi-finiteness makes $Z$ finite on an ambient
neighborhood of $L_0$. The zero relative cotangent spaces at the
six reduced fiber points give unramifiedness there. Remove the
closed image of the complement of the unramified locus to obtain
a finite unramified morphism over the whole chosen neighborhood.
This does not assert that this larger-format morphism is flat.

In the fivefold fiber product of $Z$ over that neighborhood, let $D$
be the locus of ordered pairwise distinct points. The diagonal of
a separated unramified morphism is open and closed. Each coincidence
locus is therefore open and closed, and so is their complement $D$.
It follows that $D$ is finite over the ambient neighborhood.
Every tuple in its fiber over $L_0$ is independent. Dependence is
closed on $D$; removing its closed image makes every tuple in $D$
independent after another shrinking.

The image of $D$ is closed and is now exactly the locus of five-planes
spanned by five distinct Segre points. It is contained in $\mathcal H$
and contains its dense defining span locus restricted to this
neighborhood. Since closure commutes with restriction to an open set,
it equals $\mathcal H$ there. On the graph chart all possible rank-one
points lift the six labeled core points, and every choice of five has
independent core projections. Thus this image is precisely
$\bigcup_i B_i$.
The \'etale pullback of the reduced variety $\mathcal H$ is reduced,
and every $B_i$ is smooth and reduced. Equality of these closed sets
is therefore equality of the reduced schemes. In particular, coincident
or dependent limiting configurations create no additional branch.

Eliminating $K'$ in~\eqref{nl:straight-branches}, the local span germ
is the union of six coordinate $q$-planes over the smooth base
$(L,P,Q)$ of dimension $20+3\alpha+3\beta$. Its ideal is generated
by all products of coordinates from different residue groups:
a monomial belongs to the intersection of the six branch ideals
exactly when its support meets at least two groups. The same statement
holds after completion, either by this monomial description or by
flatness of completion applied to the finite intersection of the
branch ideals.

\medskip
\emph{From five-planes to tensors.}
The flattening $C^*\to A\otimes B$ of $T$ has rank five and image
$L_0$. On the open rank-five locus, $X$ is the bundle of surjections
$C^*\twoheadrightarrow L$ over $\mathcal H$.
Indeed, for a span of five independent rank-one matrices, every
surjection onto that span gives a five-term tensor decomposition.
Conversely, a rank-five tensor with a five-term decomposition has
independent first-two-factor matrices and independent third factors.
Such tensors form a dense subset. Taking closures commutes with the
smooth open bundle projection, so the reduced rank-five locus is
exactly the indicated bundle.

Its fiber is an open subset of $\Hom(C^*,L)$, of dimension $5c$,
and it is locally trivial near $T$. It therefore adds $5c$ formal
parameters to the span germ. The $20$ smooth core parameters, the
$3\alpha+3\beta$ kernel parameters and these $5c$ tensor parameters
give
\[
 h=20+3\alpha+3\beta+5c=3a+3b+5c+2.
\]
The \'etale base change preserves the completed local ring at the
chosen complex point. Together with the span-germ calculation this
proves~\eqref{nl:completed-ring}.

\medskip
\emph{Depth and numerical invariants.}
Put
\[
 R_q=\C[[x_{i,j}:1\leq i\leq6,\ 1\leq j\leq q]]
                 /(x_{i,j}x_{k,\ell}:i\neq k),\qquad
 B_q=\prod_{i=1}^6\C[[x_{i,1},\ldots,x_{i,q}]].
\]
Each element of $R_q$ has a unique expression as a constant plus
six series, each of zero constant term and supported in a single
residue group. Hence $R_q$ embeds in $B_q$ as the tuples with a
common constant term, and there is an exact sequence of $R_q$-modules
\begin{equation}\label{nl:depth-sequence}
 0\longrightarrow R_q\longrightarrow B_q
        \longrightarrow\C^5\longrightarrow0.
\end{equation}
The last map subtracts the sixth constant term from each of the first
five. The ring $R_q$ is reduced of dimension $q$. Every branch factor
of $B_q$ is a regular $q$-dimensional quotient of $R_q$, so its
depth as an $R_q$-module is $q$. Depth of a finite direct sum is the
minimum of the depths of the summands, giving $\depth_{R_q} B_q=q$.
The nonzero finite-length module $\C^5$ has depth zero. Our range
gives $q=2(a+b-6)\geq2$; the depth lemma applied to
\eqref{nl:depth-sequence} therefore gives
\[
 \depth R_q=1.
\]

Multiplication by a formal variable is injective over every
coefficient ring, and successive quotients remove one variable.
Thus $z_1,\ldots,z_h$ is a regular sequence in
$R_q[[z_1,\ldots,z_h]]$. Its depth is $h+1$ and its dimension is
$h+q$. Dimension and depth are preserved by completion of a
Noetherian local ring, proving~\eqref{nl:dimension}
and~\eqref{nl:depth}.
The defining equations in~\eqref{nl:completed-ring} are quadratic,
so the embedding dimension is $h+6q$.
The associated graded Hilbert series is
\begin{equation}\label{nl:hilbert-series}
 (1-t)^{-h}\bigl(6(1-t)^{-q}-5\bigr).
\end{equation}
Multiplication by $(1-t)^{h+q}$ and evaluation at $t=1$ gives
multiplicity six. Embedding dimension and Hilbert--Samuel
multiplicity also agree with those of the original local ring,
which proves~\eqref{nl:embedding-multiplicity}.

\medskip
\emph{Normalization, conductor and the common intersection.}
The normalization of the completed ring is the product of its
six regular branch rings in~\eqref{nl:normalization}.
In that product, the original ring consists of tuples having the
same restriction to $\C[[z_1,\ldots,z_h]]$.
If such a tuple belongs to the conductor, multiplying by the
idempotent of any one branch must again yield a tuple with common
restriction. The other five restrictions are zero, so the original
common restriction must be zero. Conversely, a tuple whose six
restrictions are zero remains of common zero restriction after
multiplication by any element of the product. It therefore belongs
to the conductor.

This identifies the conductor with the ideal generated by all
$x_{i,j}$. Its quotient is $\C[[z_1,\ldots,z_h]]$, so the ideal is
radical. The sum of any two distinct branch ideals also kills all
the $x$ variables. Every pair of branches therefore meets in
exactly this same reduced smooth $h$-dimensional formal germ.
Finally, a simultaneous permutation of the tensor factors is a
linear isomorphism of the ambient tensor spaces and carries the
specified witness and its local ring to the transported ones.
\end{proof}

\begin{remark}\label{nl:scope}
For the padded $6\times6\times6$ witness, $h=68$ and $q=12$,
so the dimension is $80$ and the depth is $69$.
For the $3\times4\times5$ witness, $h=48$ and $q=2$,
so the dimension is $50$ and the depth is $49$.
In general the difference between dimension and depth is $q-1$.
The common branch locus has codimension $q$ in this local model
and supports its $S_2$ defect. These statements concern the explicit
rank-five witnesses. They give no formula for the depth at the
homogeneous vertex, no global normalization description, and no
higher-rank fiber formula. In particular, depth $69$ at the
$6\times6\times6$ witness is not a claim of depth less than two there.
\end{remark}
